
\documentclass{article}
\usepackage{colortbl}
\usepackage{makecell}
\usepackage{multirow}
\usepackage{supertabular}

\begin{document}

\newcounter{utterance}

\centering \large Interaction Transcript for game `dond', experiment `semi\_en', episode 11 with Aya{-}Expanse{-}32B{-}t0.0.
\vspace{24pt}

{ \footnotesize  \setcounter{utterance}{1}
\setlength{\tabcolsep}{0pt}
\begin{supertabular}{c@{$\;$}|p{.15\linewidth}@{}p{.15\linewidth}p{.15\linewidth}p{.15\linewidth}p{.15\linewidth}p{.15\linewidth}}
    \# & $\;$A & \multicolumn{4}{c}{Game Master} & $\;\:$B\\
    \hline

    \theutterance \stepcounter{utterance}  
    & & \multicolumn{4}{p{0.6\linewidth}}{
        \cellcolor[rgb]{0.9,0.9,0.9}{
            \makecell[{{p{\linewidth}}}]{
                \texttt{\tiny{[P1$\langle$GM]}}
                \texttt{You are playing a negotiation game in which you have to agree on how to divide a set of items among you and another player.} \\
\\ 
\texttt{Rules:} \\
\texttt{(a) You and the other player are given a set of items. Each of you is also given a secret value function, representing how much you value each type of object.} \\
\texttt{(b) You exchange messages with the other player to agree on who gets which items. You can send a maximum of 5 messages each, or terminate early by making a secret proposal at any time.} \\
\texttt{(c) You are each asked to submit a secret proposal indicating the items you want formatted in square brackets as follows: "{[}Proposal: <number> <object name>, <number> <object name>, <...>{]}"} \\
\texttt{(d) If your proposals are complementary, i.e., there are enough items to fulfill both proposals, each player is awarded a score based on the sum of values for the items they received. Otherwise, both of you get zero points.} \\
\texttt{(e) Your goal is to maximize the score you receive.} \\
\\ 
\texttt{Let us start.} \\
\\ 
\texttt{The set of available items is:} \\
\\ 
\texttt{1 pot, 1 candy, 1 whistle, 2 dice} \\
\\ 
\texttt{Your secret value function is:} \\
\\ 
\texttt{pot: 2, candy: 3, whistle: 3, dice: 1} \\
\\ 
\texttt{IMPORTANT: Your messages, unless it is the secret proposal, are directly transmitted to the other player, so do not include any response to the rules or text announcing your message. To make a secret proposal, use the indicated format. Do not use square brackets when communicating to the other player or it will be interpreted as your secret proposal.} \\
            }
        }
    }
    & \\ \\

    \theutterance \stepcounter{utterance}  
    & \multicolumn{4}{p{0.6\linewidth}}{
        \cellcolor[rgb]{0.9,0.9,0.9}{
            \makecell[{{p{\linewidth}}}]{
                \texttt{\tiny{[P1$\rangle$GM]}}
                \texttt{Understanding the rules and my valuation, I'll initiate the negotiation by messaging the other player:} \\
\\ 
\texttt{**Message 1:** "Hello, I value the candy and whistle highly. Would you be interested in a trade? I can offer both of them for the 2 dice."} \\
\\ 
\texttt{After this message, I'll await their response. Depending on their reaction, I'll adjust my strategy. If they seem open to trading, I might send another message offering a specific division. If they propose something different, I'll evaluate their offer and decide whether to accept, counter, or make my secret proposal.} \\
\\ 
\texttt{**Potential Secret Proposal (to be submitted later if necessary):**} \\
\texttt{"Proposal: 1 candy, 1 whistle"} \\
\\ 
\texttt{This strategy aims to maximize my score based on my value function while adhering to the game's rules and the need for complementary proposals.} \\
            }
        }
    }
    & & \\ \\

    \theutterance \stepcounter{utterance}  
    & & & \multicolumn{2}{p{0.3\linewidth}}{
        \cellcolor[rgb]{0.9,0.9,0.9}{
            \makecell[{{p{\linewidth}}}]{
                \texttt{\tiny{[GM$|$GM]}}
                \texttt{continue} \\
            }
        }
    }
    & & \\ \\

    \theutterance \stepcounter{utterance}  
    & & & \multicolumn{2}{p{0.3\linewidth}}{
        \cellcolor[rgb]{0.9,0.9,0.9}{
            \makecell[{{p{\linewidth}}}]{
                \texttt{\tiny{[GM$|$GM]}}
                \texttt{Understanding the rules and my valuation, I'll initiate the negotiation by messaging the other player:} \\
\\ 
\texttt{**Message 1:** "Hello, I value the candy and whistle highly. Would you be interested in a trade? I can offer both of them for the 2 dice."} \\
\\ 
\texttt{After this message, I'll await their response. Depending on their reaction, I'll adjust my strategy. If they seem open to trading, I might send another message offering a specific division. If they propose something different, I'll evaluate their offer and decide whether to accept, counter, or make my secret proposal.} \\
\\ 
\texttt{**Potential Secret Proposal (to be submitted later if necessary):**} \\
\texttt{"Proposal: 1 candy, 1 whistle"} \\
\\ 
\texttt{This strategy aims to maximize my score based on my value function while adhering to the game's rules and the need for complementary proposals.} \\
            }
        }
    }
    & & \\ \\

    \theutterance \stepcounter{utterance}  
    & & \multicolumn{4}{p{0.6\linewidth}}{
        \cellcolor[rgb]{0.9,0.9,0.9}{
            \makecell[{{p{\linewidth}}}]{
                \texttt{\tiny{[GM$\rangle$P2]}}
                \texttt{You are playing a negotiation game in which you have to agree on how to divide a set of items among you and another player.} \\
\\ 
\texttt{Rules:} \\
\texttt{(a) You and the other player are given a set of items. Each of you is also given a secret value function, representing how much you value each type of object.} \\
\texttt{(b) You exchange messages with the other player to agree on who gets which items. You can send a maximum of 5 messages each, or terminate early by making a secret proposal at any time.} \\
\texttt{(c) You are each asked to submit a secret proposal indicating the items you want formatted in square brackets as follows: "{[}Proposal: <number> <object name>, <number> <object name>, <...>{]}"} \\
\texttt{(d) If your proposals are complementary, i.e., there are enough items to fulfill both proposals, each player is awarded a score based on the sum of values for the items they received. Otherwise, both of you get zero points.} \\
\texttt{(e) Your goal is to maximize the score you receive.} \\
\\ 
\texttt{Let us start.} \\
\\ 
\texttt{The set of available items is:} \\
\\ 
\texttt{1 pot, 1 candy, 1 whistle, 2 dice} \\
\\ 
\texttt{Your secret value function is:} \\
\\ 
\texttt{pot: 8, candy: 0, whistle: 0, dice: 1} \\
\\ 
\texttt{IMPORTANT: Your messages, unless it is the secret proposal, are directly transmitted to the other player, so do not include any response to the rules or text announcing your message. To make a secret proposal, use the indicated format. Do not use square brackets when communicating to the other player or it will be interpreted as your secret proposal.} \\
\\ 
\\ 
\texttt{Understanding the rules and my valuation, I'll initiate the negotiation by messaging the other player:} \\
\\ 
\texttt{**Message 1:** "Hello, I value the candy and whistle highly. Would you be interested in a trade? I can offer both of them for the 2 dice."} \\
\\ 
\texttt{After this message, I'll await their response. Depending on their reaction, I'll adjust my strategy. If they seem open to trading, I might send another message offering a specific division. If they propose something different, I'll evaluate their offer and decide whether to accept, counter, or make my secret proposal.} \\
\\ 
\texttt{**Potential Secret Proposal (to be submitted later if necessary):**} \\
\texttt{"Proposal: 1 candy, 1 whistle"} \\
\\ 
\texttt{This strategy aims to maximize my score based on my value function while adhering to the game's rules and the need for complementary proposals.} \\
            }
        }
    }
    & \\ \\

    \theutterance \stepcounter{utterance}  
    & & & \multicolumn{4}{p{0.6\linewidth}}{
        \cellcolor[rgb]{0.9,0.9,0.9}{
            \makecell[{{p{\linewidth}}}]{
                \texttt{\tiny{[GM$\langle$P2]}}
                \texttt{**Message 2:** "I understand your interest in the candy and whistle. However, I value the pot the most. Perhaps we could reach a compromise? I'm open to discussing a fair trade that benefits us both."} \\
\\ 
\texttt{I want to emphasize the importance of a mutually beneficial agreement, given the game's rules about complementary proposals. This approach might encourage the other player to consider a division that includes the pot, which they seem to value more than the dice.} \\
\\ 
\texttt{I'll await their response and be prepared to discuss specific items or consider their proposal. If necessary, I can make a counteroffer or submit my secret proposal to ensure I don't exceed the 5{-}message limit.} \\
\\ 
\texttt{**Potential Next Steps:**} \\
\texttt{{-} If they express interest in the pot, I might propose: "How about I take the pot and one die, and you get the rest?"} \\
\texttt{{-} If they seem unwilling to part with the pot, I can submit my secret proposal, knowing that my score will depend on the outcome of their proposal.} \\
            }
        }
    }
     \\ \\

    \theutterance \stepcounter{utterance}  
    & & & \multicolumn{2}{p{0.3\linewidth}}{
        \cellcolor[rgb]{0.9,0.9,0.9}{
            \makecell[{{p{\linewidth}}}]{
                \texttt{\tiny{[GM$|$GM]}}
                \texttt{continue} \\
            }
        }
    }
    & & \\ \\

    \theutterance \stepcounter{utterance}  
    & & & \multicolumn{2}{p{0.3\linewidth}}{
        \cellcolor[rgb]{0.9,0.9,0.9}{
            \makecell[{{p{\linewidth}}}]{
                \texttt{\tiny{[GM$|$GM]}}
                \texttt{**Message 2:** "I understand your interest in the candy and whistle. However, I value the pot the most. Perhaps we could reach a compromise? I'm open to discussing a fair trade that benefits us both."} \\
\\ 
\texttt{I want to emphasize the importance of a mutually beneficial agreement, given the game's rules about complementary proposals. This approach might encourage the other player to consider a division that includes the pot, which they seem to value more than the dice.} \\
\\ 
\texttt{I'll await their response and be prepared to discuss specific items or consider their proposal. If necessary, I can make a counteroffer or submit my secret proposal to ensure I don't exceed the 5{-}message limit.} \\
\\ 
\texttt{**Potential Next Steps:**} \\
\texttt{{-} If they express interest in the pot, I might propose: "How about I take the pot and one die, and you get the rest?"} \\
\texttt{{-} If they seem unwilling to part with the pot, I can submit my secret proposal, knowing that my score will depend on the outcome of their proposal.} \\
            }
        }
    }
    & & \\ \\

    \theutterance \stepcounter{utterance}  
    & & \multicolumn{4}{p{0.6\linewidth}}{
        \cellcolor[rgb]{0.9,0.9,0.9}{
            \makecell[{{p{\linewidth}}}]{
                \texttt{\tiny{[P1$\langle$GM]}}
                \texttt{**Message 2:** "I understand your interest in the candy and whistle. However, I value the pot the most. Perhaps we could reach a compromise? I'm open to discussing a fair trade that benefits us both."} \\
\\ 
\texttt{I want to emphasize the importance of a mutually beneficial agreement, given the game's rules about complementary proposals. This approach might encourage the other player to consider a division that includes the pot, which they seem to value more than the dice.} \\
\\ 
\texttt{I'll await their response and be prepared to discuss specific items or consider their proposal. If necessary, I can make a counteroffer or submit my secret proposal to ensure I don't exceed the 5{-}message limit.} \\
\\ 
\texttt{**Potential Next Steps:**} \\
\texttt{{-} If they express interest in the pot, I might propose: "How about I take the pot and one die, and you get the rest?"} \\
\texttt{{-} If they seem unwilling to part with the pot, I can submit my secret proposal, knowing that my score will depend on the outcome of their proposal.} \\
            }
        }
    }
    & \\ \\

    \theutterance \stepcounter{utterance}  
    & \multicolumn{4}{p{0.6\linewidth}}{
        \cellcolor[rgb]{0.9,0.9,0.9}{
            \makecell[{{p{\linewidth}}}]{
                \texttt{\tiny{[P1$\rangle$GM]}}
                \texttt{**Message 3:** "I appreciate your willingness to compromise. Given our value functions, it seems the pot is a key item for both of us. How about this: I'll take the pot and one die, and you can have the candy, whistle, and the other die."} \\
\\ 
\texttt{This proposal aims to find a middle ground, ensuring that both players receive items they value highly. By offering the pot and one die, I'm hoping to secure a positive response, as it directly addresses their interest in the pot.} \\
\\ 
\texttt{**Potential Next Steps:**} \\
\texttt{{-} If they accept, we can finalize the agreement, and I'll submit my secret proposal accordingly.} \\
\texttt{{-} If they counter with a different division, I'll evaluate it based on my value function and the game's rules. I might say, "Let's adjust it slightly. How about I take the pot and whistle, and you get the candy, dice, and an extra point for fairness?"} \\
\texttt{{-} If negotiations stall, I can submit my secret proposal, "{[}Proposal: 1 pot, 1 whistle{]}," as it's a fair offer based on our values.} \\
            }
        }
    }
    & & \\ \\

    \theutterance \stepcounter{utterance}  
    & & & \multicolumn{2}{p{0.3\linewidth}}{
        \cellcolor[rgb]{0.9,0.9,0.9}{
            \makecell[{{p{\linewidth}}}]{
                \texttt{\tiny{[GM$|$GM]}}
                \texttt{proposal} \\
            }
        }
    }
    & & \\ \\

    \theutterance \stepcounter{utterance}  
    & & & \multicolumn{2}{p{0.3\linewidth}}{
        \cellcolor[rgb]{0.9,0.9,0.9}{
            \makecell[{{p{\linewidth}}}]{
                \texttt{\tiny{[GM$|$GM]}}
                [1, 0, 1, 0]
            }
        }
    }
    & & \\ \\

    \theutterance \stepcounter{utterance}  
    & & \multicolumn{4}{p{0.6\linewidth}}{
        \cellcolor[rgb]{0.9,0.9,0.9}{
            \makecell[{{p{\linewidth}}}]{
                \texttt{\tiny{[GM$\rangle$P2]}}
                \texttt{The other player has submitted a secret proposal. It is now your turn to submit a proposal of your own.} \\
            }
        }
    }
    & \\ \\

    \theutterance \stepcounter{utterance}  
    & & & \multicolumn{4}{p{0.6\linewidth}}{
        \cellcolor[rgb]{0.9,0.9,0.9}{
            \makecell[{{p{\linewidth}}}]{
                \texttt{\tiny{[GM$\langle$P2]}}
                \texttt{Given the other player's secret proposal and my understanding of their value function, I submit the following secret proposal:} \\
\\ 
\texttt{**Secret Proposal:** "Proposal: 1 pot, 1 die"} \\
\\ 
\texttt{This proposal aims to secure the item I value the most (the pot) while offering a fair division of the remaining items. By taking one die, I leave the other die and the candy for the other player, which should be an attractive offer based on their value function.} \\
\\ 
\texttt{This strategy maximizes my score while considering the other player's interests, increasing the likelihood of complementary proposals and a successful negotiation.} \\
            }
        }
    }
     \\ \\

    \theutterance \stepcounter{utterance}  
    & & & \multicolumn{2}{p{0.3\linewidth}}{
        \cellcolor[rgb]{0.9,0.9,0.9}{
            \makecell[{{p{\linewidth}}}]{
                \texttt{\tiny{[GM$|$GM]}}
                \texttt{continue} \\
            }
        }
    }
    & & \\ \\

    \theutterance \stepcounter{utterance}  
    & & & \multicolumn{2}{p{0.3\linewidth}}{
        \cellcolor[rgb]{0.9,0.9,0.9}{
            \makecell[{{p{\linewidth}}}]{
                \texttt{\tiny{[GM$|$GM]}}
                \texttt{player was instructed to make a proposal} \\
            }
        }
    }
    & & \\ \\

\end{supertabular}
}

\end{document}
